\statement{B.109}{Proved}{The three natural mixed-neck maps have actual node kernels, with two distinct marked relations at the first mixed neck}
\label{stmt:B109}
Use the actual surface $\mathscr Q$ and involution of B.106, and put
\[
 Q=\frac{1+q}{1-q},\qquad Z=\frac{t+1}{t-1},\qquad
 c_j=\cot\frac{j\pi}{5},\qquad \lambda_j=\frac{2}{1+c_j^2}
 \quad(1\le j\le4).
\]
The inner and outer bases are respectively $\Re Q>0$ and $\Re Q<0$,
with their actual old punctures retained. The moving punctures are
\begin{equation}\label{eq:B109-points}
 A_j=\lambda_j(Q-i c_j),\qquad B_j=A_j^{-1},\qquad
 Z^2-2QZ+1=0,
\end{equation}
and the two retained fixed punctures are $Z=0,\infty$.
The involution is $Z\mapsto Z^{-1}$.

For each $k=1,2,3$, use the natural lower-$t$-semicircle marking on
the unit-circle neck between $e^{2\pi i k/5}$ and
$e^{2\pi i(k+1)/5}$. Its selected upper-$Z$ punctures are
\[
 B_1,\ldots,B_k,A_{k+1},\ldots,A_4
 \quad\hbox{and one retained branch puncture}.
\]
They define a geometric free group $K_k\cong F_5$. With the same
regular-neck transport on the two sides, the natural maps from
$K_k$ into \emph{both} actual bad regional orbifold vertex groups
are noninjective. Each kernel contains a specified transported
ordinary-node commutator. This assertion concerns these natural
markings, not every conceivable common rank-five subgroup.

For the first mixed neck $k=1$, there is one explicitly based
geometric basis $x_1,\ldots,x_5$, in physical order
\[
 A_2,A_4,A_3,\text{upper branch},B_1,
\]
in which the two actual node relations have normal-closure
representatives
\begin{equation}\label{eq:B109-two-relations}
 [x_1,x_5],\qquad [x_1,w x_5w^{-1}],\qquad w=x_2x_3x_4.
\end{equation}
The first comes from the outer side and the second from the inner
side. The quotient of $K_1$ by \emph{precisely these two relations}
is
\begin{equation}\label{eq:B109-block}
 \langle a,b,w\mid[a,b],[a,wbw^{-1}]\rangle
       *F(x_3,x_4),\qquad a=x_1, b=x_5.
\end{equation}
Its first factor is the HNN extension of
$\mathbb Z(a)\times F(b,d)$ with relation $wbw^{-1}=d$.
Its kernel under $w\mapsto1$, $a,b\mapsto0$, to $\mathbb Z$ has
the presentation
\begin{equation}\label{eq:B109-path}
 \left\langle a_n,b_n\ (n\in\mathbb Z)\;\middle|\;
 [a_n,b_n],\ [a_n,b_{n+1}]\ (n\in\mathbb Z)\right\rangle,
 \quad a_n=w^naw^{-n},\ b_n=w^nbw^{-n}.
\end{equation}
No presentation of a full regional image or of the global surface
group is asserted by \eqref{eq:B109-block}.

\begin{proof}
\emph{B.109.1. The actual coordinates and the three markings.}
The identities in \eqref{eq:B109-points} follow by substituting
B.106's $T_u$ with $u=e^{2\pi i j/5}$, using
$1-\cos(2\pi j/5)=\lambda_j$ and
$\sin(2\pi j/5)=\lambda_jc_j$. The branch equation follows from
$qt^2=1$. In the strip
\begin{equation}\label{eq:B109-strip}
 c_{k+1}<\Im Q<c_k,
\end{equation}
the imaginary signs of $A_j$ and $B_j$ are opposite, and the
positive ones are exactly the four stated labels. Moreover
\[
 \Im t=-\frac{2\Im Z}{|Z-1|^2},
\]
so this is the natural lower-$t$-semicircle lift, not an arbitrary
choice of one peripheral from each pair.

The branch roots have product one. If $Q$ is nonreal, neither root
is real, and their imaginary parts have opposite signs. If $Q$ is
real with $|Q|<1$, the two roots are nonreal conjugates, again with
opposite imaginary signs. For $k=1,3$ the strip avoids the real
axis; for $k=2$ below we restrict to $|\Re Q|<1$. Thus the selected
branch is defined throughout each path. All ten moving points have
real part with the same sign as $\Re Q$: this is immediate for
$A_j,B_j$, while a branch root satisfies
\[
 \Re Q=\frac{\Re Z}{2}(1+|Z|^{-2}).
\]

In a boundary fiber take a sufficiently large positive imaginary
$Z$ basepoint above all five selected points, with one connector
in the selected half-plane. The whole punctured half-plane gives
$K_k=F_5$. An annular-sector disk in $t^2$, cut at the retained
point $1$, has this same lift and group. A single path from the
literal point $\bar P=(399/40000,0)$ of B.106 to this fiber
basepoint identifies these groups with based subgroups of the
actual surface and its orbifold extension. That path is fixed for
both regional maps. Changing it conjugates the whole diagram and
does not change an injection or kernel assertion. No comparison
with an earlier Artin array is silently made.

\emph{B.109.2. Complete critical coordinates along the chosen paths.}
Put
\[
 c=\sqrt{1+2/\sqrt5},\qquad d=\sqrt{1-2/\sqrt5};
 \qquad(c_1,c_2,c_3,c_4)=(c,d,-d,-c),\quad cd=1/\sqrt5.
\]
For $i\ne j$, the equation $A_i=B_j$ is
\[
 \lambda_i\lambda_j(Q-i c_i)(Q-i c_j)=1.
\]
Its two roots are
\begin{equation}\label{eq:B109-critical}
 Q_{ij}^{\pm}=\pm\frac{|1+c_ic_j|}{2}
                    +\frac{i(c_i+c_j)}2.
\end{equation}
Indeed the imaginary equation gives $\Im Q=(c_i+c_j)/2$, and
\[
 \frac1{\lambda_i\lambda_j}-\frac{(c_i-c_j)^2}{4}
       =\frac{(1+c_ic_j)^2}{4}.
\]
The same formula for $i=j$ gives the two cone-contact values.
The positive real part and the imaginary part of all ten inner
unordered-pair values are the following:
\[
\begin{array}{c|c|c}
 ij&\Re Q_{ij}^{+}&\Im Q_{ij}^{+}\\\hline
 11&(1+c^2)/2&c\\
 12&(1+cd)/2&(c+d)/2\\
 13&(1-cd)/2&(c-d)/2\\
 14&(c^2-1)/2&0\\
 22&(1+d^2)/2&d\\
 23&(1-d^2)/2&0\\
 24&(1-cd)/2&-(c-d)/2\\
 33&(1+d^2)/2&-d\\
 34&(1+cd)/2&-(c+d)/2\\
 44&(1+c^2)/2&-c
\end{array}
\]
The two zero-height entries coincide at $1/\sqrt5$. The outer
table changes only the sign of the real part. B.106's complete
factorization puts every same-sign or old-branch collision on
the removed unit-circle boundary; the remaining old interior
point $q=0$ is $Q=1$, and the omitted point $q=\infty$ is $Q=-1$.

Set
\[
 x_*=(1+1/\sqrt5)/2,\qquad y_*=(c+d)/2.
\]
For the three middle necks use respectively
\[
 Q_{*,1}=iy_*,\quad Q_{*,2}=0,\quad Q_{*,3}=-iy_*.
\]
These are regular unit-circle points in the indicated arcs.
The actual horizontal paths from them end at
\[
 \pm x_*+iy_*,\qquad \pm1/\sqrt5,\qquad \pm x_*-iy_*.
\]
The table proves that no open path meets an earlier critical
value. For $k=1,3$ the endpoint has only the pair $(1,2)$ or
$(3,4)$, together with its involution mate. For $k=2$ both pairs
$(1,4)$ and $(2,3)$ occur at the endpoint, in two disjoint upper
collision disks and their two lower mates. Both nodes are kept.
Also $1/\sqrt5<1$, so the middle paths avoid $Q=\pm1$ and keep
the branch roots nonreal. Each crossing is transverse: vanishing
of the derivative difference
$\lambda_i+\lambda_j/A_j^2$ would make $A_j$ purely imaginary,
contrary to its nonzero real part at the target.

\emph{B.109.3. Actual based vanishing disks and noninjection.}
First move an arbitrarily short distance through the regular neck
to $Q_{*,k}+\epsilon$ on the inner side, or
$Q_{*,k}-\epsilon$ on the outer side. Transport the common
half-plane marking and connector through that same neck. At these
endpoints the selected disk can be isotoped into the corresponding
upper-right or upper-left quadrant. No such assertion is made for
the boundary fiber itself, where the punctures have zero real part.

Continue along the displayed horizontal path. The selected points
remain in that quadrant, with no collision before its endpoint.
In a small upper collision disk take the actual local vanishing
arc between $A_{k+1}$ and $B_k$, and transport it and its meridian
whiskers backward in this literal punctured-half-plane bundle.
They give a pair in a geometric free basis of the same $K_k$.
Their commutator is nontrivial in $K_k$. Inserting the actual
ordinary-node critical disk makes that commutator trivial, by
the local normal-crossing bidisk complement. This nullhomotopy
already lies in the bad quadrant and hence in the bad orbifold
vertex. For $k=2$ choose either of the two disjoint upper node
disks; its simultaneous companion changes neither assertion.

This proves noninjection on both sides. The two transported words
are not thereby identified: the two paths can pass on different
sides of the other selected punctures. Transport along the regular
unit-circle arc gives the same conclusion for any other basepoint
of the same natural middle-neck marking.

\emph{B.109.4. Explicit vanishing arcs at the first mixed neck.}
For $k=1$ put $y=y_*$, $\delta=(c-d)/2>0$ and $Q=x+iy$.
The two selected endpoints are
\[
 A_2=\lambda_2(x+i\delta),\qquad
 B_1=\frac{x+i\delta}{\lambda_1(x^2+\delta^2)}.
\]
For $|x|<x_*$ they are distinct positive multiples of the same
vector. Their straight radial segment is therefore a literal
shrinking arc. For $j=3,4$ its cross product with $A_j$ is
$\lambda_jx(d-c_j)$, which has the sign of $x$ and never vanishes
away from the boundary fiber.

The retained branch is controlled as well. Write its upper root
as $X+iY$, $r^2=|Z|^2$. Then
\[
 x=X(1+r^{-2})/2,\qquad y=Y(1-r^{-2})/2.
\]
Since $y,Y>0$, we have $r>1$. Its cross product with
$(x,\delta)$ is
\[
 X\bigl(Y(1+r^{-2})/2-\delta\bigr).
\]
The bracket exceeds $y-\delta=d>0$, so it too has the sign of
$x$. Thus all three middle punctures stay strictly on the same
side of the radial arc. No uncomputed threading around the branch
has been removed.

At $x=0$ the five upper heights, in order, are
\[
 \lambda_2\delta,\quad \lambda_1(y+c),\quad
 \lambda_2(y+d),\quad y+\sqrt{1+y^2},\quad
 \frac1{\lambda_1\delta}.
\]
They lie respectively in
\[
 (0.95,0.96),\ (1.53,1.55),\ (2.12,2.14),\
 (2.16,2.17),\ (2.75,2.76).
\]
These are exact rational inequalities: substitute the displayed
positive radicals and square the positive sides. For example the
bounds $2.23606<\sqrt5<2.23607$,
$1.37638<c<1.37639$ and $0.32491<d<0.32492$, each checked by
squaring, give these wider intervals by monotonic interval
arithmetic. Therefore the inner arc approaches the right side
of all three middle points in the $Z$ plane; the outer arc is
its left-side reflection. A sufficiently short regular transport
to the boundary moves each puncture in a disjoint small disk,
so these are the actual two arc isotopy classes there.

\emph{B.109.5. One calibrated basis and both raw words.}
Rotate by $z=-iZ$. The five punctures are now on the positive real
axis in the order $A_2,A_4,A_3,\text{branch},B_1$. An increasing
real-axis homeomorphism, extended preserving both half-planes,
moves them to $1,2,3,4,5$. Use basepoint $6$ and downward cuts;
a positive meridian $x_j$ crosses its cut from left to right.
The common connector is
\[
 6\longrightarrow6+i\longrightarrow\tfrac12+i
      \longrightarrow p=\tfrac12.
\]
It has empty cut word and stays away from the retained boundary
point. The first meridian is $a=x_1$. A path from $p$ above the
points to a positive meridian at $5$ gives $b_{\rm out}=x_5$.
A path below them crosses the first four downward cuts positively
before reaching the same small positive meridian, and gives
\[
 b_{\rm in}=(x_1x_2x_3x_4)x_5(x_1x_2x_3x_4)^{-1}.
\]
The thin neighborhoods of these two arc trees supply the actual
local commutator disks. Thus the raw words are $[x_1,x_5]$ and
$[x_1,b_{\rm in}]$. Conjugating the \emph{whole} second word by
$x_1^{-1}$ gives $[x_1,wx_5w^{-1}]$, proving
\eqref{eq:B109-two-relations}. This step does not discard an
individual peripheral conjugator. Reversing a node disk would
only invert its defining relation.

\emph{B.109.6. The exact auxiliary quotient and its limits.}
The replacement of $(x_2,x_3,x_4)$ by $(w,x_3,x_4)$ is an actual
free-basis change, with $x_2=wx_4^{-1}x_3^{-1}$. This proves
\eqref{eq:B109-block}. Introducing $d=wbw^{-1}$ gives
\[
 \langle a,b,d,w\mid[a,b],[a,d],wbw^{-1}=d\rangle.
\]
The base is $\mathbb Z(a)\times F(b,d)$, and the two associated
infinite cyclic subgroups inject in it. Thus this is a genuine
HNN extension. Reidemeister--Schreier rewriting with representatives
$w^n$ gives \eqref{eq:B109-path}; the second commutator has index
$n+1$, in accordance with the specified conjugation convention.
The unused free factors contribute their free $w$-translates if
one takes the kernel for the whole group in
\eqref{eq:B109-block}.

The exact coordinate table and horizontal-path checks were
executed in
\texttt{work\_DE/next\_marginal\_source/check\_middle\_neck\_nodes21.py};
the initial symbolic-expression comparison failure and the
corrected genuine output are both retained. The separate
\texttt{check\_middle\_neck\_arc\_words21.py} checks every rational
polygon segment, cut crossing, positive small-loop orientation,
raw word and whole-relator conjugation. CE's
\texttt{check\_first\_mixed\_neck\_arc.py} independently supplies
the radical intervals, and both independent full-proof reviews
retain the branch-side and initial-collar arguments.

Only the natural middle-neck injections have been refuted.
The two-relator quotient does not contain all regional relations,
all old base operators, or the full global node orbit. Consequently
no global fixed-source failure, global amalgam, full kernel
freeness, or vanishing of the residual obstruction follows.
\end{proof}
